\section{The mathematical definition}\label{sec:06-groups:01-definition:1}

A \textbf{group} is a set \(G\) together with:

\begin{itemize}
\tightlist
\item
  a binary operation \(\cdot : G \times G \to G\),
\item
  a distinguished element \(e \in G\) (the identity),
\item
  an inverse function \((-)^{-1} : G \to G\),
\end{itemize}

satisfying three axioms, for all \(a, b, c \in G\):

\[
\begin{aligned}
\text{(associativity)} &\quad (a \cdot b) \cdot c = a \cdot (b \cdot c) \\
\text{(identity)} &\quad e \cdot a = a \quad\text{and}\quad a \cdot e = a \\
\text{(inverse)} &\quad a^{-1} \cdot a = e \quad\text{and}\quad a \cdot a^{-1} = e
\end{aligned}
\]

No further axiom is required; in particular, commutativity is not assumed in general (a group where \(a \cdot b = b \cdot a\) always holds is called \textbf{abelian} or \textbf{commutative}, defined separately below).

\subsection{References}\label{sec:06-groups:01-definition:2}

Full citations in the \hyperref[sec:root:bibliography:1]{Bibliography}.

\begin{itemize}
\tightlist
\item
  Dummit and Foote (\cite{DummitFoote2003}) --- the standard classical (non-Lean) reference for the group axioms above and their consequences.
\item
  Aluffi (\cite{Aluffi2009}) --- a group-theory treatment using the same categorical framing (forgetful functors, universal properties) this book uses from Chapter 6 §6 onward.
\end{itemize}
